% !TEX root = ../notes_missingsemester.tex
% Chapter 8: Automate Everything: CI/CD, Cron, Webhooks & Backups
%%%%%%%%%%%%%%%%%%%%%%%%%%%%%%%%%%%%%%%%%%%%%%%%%%%%%%%%%%%%%%%%%%%%%%

\index{CI/CD}\index{GitHub Actions}\index{cron}\index{backup}\index{feature flag}
\chapter{Automate Everything: CI/CD, Cron, Webhooks \& Backups}
\printchapterversion{08_cicd_automation}
\newthought{Manual deploys are technical debt. Automate or regret.}

% ==============================================================================================
% BLOCK 1: INTRODUCTION
% ==============================================================================================

\section{CI/CD: What It Is and Why It Matters}\index{CI/CD}

\newthought{Continuous Integration\index{continuous integration}} means every push is automatically linted and
tested. Broken code is caught before it reaches the server.

\newthought{Continuous Deployment\index{continuous deployment}} means every merge to \texttt{main} automatically
deploys to the server, with no manual SSH-and-pull.

\newthought{The pipeline:} lint $\to$ test $\to$ deploy. Each step gates the next. If linting fails,
tests do not run. If tests fail, the deploy does not happen.

% ==============================================================================================
% BLOCK 2: CONCEPTS AND PRACTICE
% ==============================================================================================

\newpage
\section{GitHub Actions}\index{GitHub Actions}

\newthought{GitHub Actions}\marginnote{GitHub Actions\index{GitHub Actions}:
\url{https://docs.github.com/en/actions}\\
Workflow syntax: \url{https://docs.github.com/en/actions/using-workflows/workflow-syntax-for-github-actions}}
run workflows in response to repository events. The configuration lives in
YAML files under \texttt{.github/workflows/}.

\begin{verbatim}
# .github/workflows/ci.yml
name: CI
on: [push, pull_request]

jobs:
  lint-and-test:
    runs-on: ubuntu-latest
    steps:
      - uses: actions/checkout@v4

      - name: Set up Python
        uses: actions/setup-python@v5
        with:
          python-version: "3.11"

      - name: Install dependencies
        run: |
          python -m venv .venv
          source .venv/bin/activate
          pip install -r requirements.txt
          pip install ruff pytest

      - name: Lint
        run: ruff check .

      - name: Test
        run: pytest tests/
\end{verbatim}

\marginnote{Key concepts: \emph{triggers}\index{GitHub Actions!triggers} (\texttt{push}, \texttt{pull\_request})
define when a workflow runs. \emph{Jobs} run on \emph{runners}, GitHub-provided VMs. \emph{Steps} are
sequential commands within a job.}

\subsection{Using Secrets in Actions}\index{GitHub Actions!secrets}

\newthought{Never hardcode secrets in workflow files.} Store them in GitHub's repository settings under
Settings $\to$ Secrets and variables $\to$ Actions. Reference them as \texttt{\$\{\{ secrets.MY\_SECRET \}\}}.
Never echo them in logs.

\subsection{The Deploy Workflow}

\begin{verbatim}
# .github/workflows/deploy.yml
name: Deploy
on:
  push:
    branches: [main]

jobs:
  deploy:
    runs-on: ubuntu-latest
    steps:
      - name: Deploy to server
        uses: appleboy/ssh-action@v1
        with:
          host: ${{ secrets.SERVER_IP }}
          username: ${{ secrets.SERVER_USER }}
          key: ${{ secrets.SSH_PRIVATE_KEY }}
          script: |
            cd /opt/pixelwise
            bash scripts/deploy.sh
\end{verbatim}

\subsection{Tests: Where They Live}\index{testing}\index{pytest}

\newthought{The CI pipeline runs \texttt{pytest},}\marginnote{pytest\index{pytest}:
\url{https://docs.pytest.org/}\\
ruff\index{ruff} (linter): \url{https://docs.astral.sh/ruff/}} but writing tests is an entire discipline covered
in dedicated courses. Here we show where tests live (\texttt{tests/}), what a test looks like, and how
the CI pipeline runs them. A few example tests are provided in the boilerplate:

\begin{verbatim}
# tests/test_classifier.py
from app.classifier import classify_batch
import numpy as np

def test_classify_batch_shape():
    images = np.zeros((2, 28, 28), dtype=np.uint8)
    results = classify_batch(images)
    assert len(results) == 2
    assert "prediction" in results[0]
    assert "confidence" in results[0]
\end{verbatim}

The goal is to see the feedback loop: push code $\to$ tests run $\to$ green or red.


% ==============================================================================================
\newpage
\section{Optional Exercise: Catch Up to \texttt{v0.6}}\index{exercises}\index{Git!tag}\index{recovery}

\newthought{Catching up via snapshot and tag.}\index{Git!tag}\index{recovery}
This block builds on the state at the end of Block~7: Nginx installed
on \texttt{prod} as a reverse proxy in front of the FastAPI service, the
static frontend served from \texttt{/opt/pixelwise/frontend/}, a
self-signed TLS certificate in place, the site reachable at
\texttt{https://192.168.56.11}, and security headers set on every
response.
If you joined late, missed Block~7, or your repository has drifted,
restore from the snapshot and reset to \texttt{v0.6} rather than redoing
every prior exercise.

\marginnote{\newthought{Tag map.}\index{Git!tag}
\texttt{v0.6} marks the end of Block~7, \texttt{v0.7} will mark the end
of this block.
Browse the full set at
\url{https://github.com/schutera/pixelwise/tags}.}

\newthought{Restore the snapshot.}
On both VMs, restore the \texttt{B7-complete} snapshot from Block~7.
The snapshot carries everything that is not in Git, the \texttt{.venv/},
the model artefact, the systemd service, the PostgreSQL installation,
the Nginx configuration, the TLS certificate, and any real \texttt{.env}
values you wrote.
With the system layer back in place, only the repository state needs
to be rewound.

\newthought{Catch-up procedure on \texttt{dev}.}
Reset PixelWise to the previous tag:

\begin{verbatim}
cd ~/pixelwise
git fetch origin
git reset --hard v0.6
\end{verbatim}

\marginnote{\newthought{Why \texttt{reset -{}-hard}?}\index{Git!reset}
This rewrites your working tree to match the end of Block~7 exactly,
including \texttt{frontend/}, the Nginx configuration committed to the
repository, and the updated CORS middleware in \texttt{app/main.py}.
There is no work to lose if you have not committed beyond \texttt{v0.6};
if you have, branch off first with \texttt{git branch backup-pre-catchup}
so the work is recoverable.}

\newthought{Mirror it on \texttt{prod}.}

\begin{verbatim}
ssh produser@192.168.56.11
cd /opt/pixelwise
git fetch origin
git reset --hard v0.6
sudo systemctl reload nginx
sudo systemctl restart pixelwise
\end{verbatim}

\marginnote{\newthought{Already have a fork?}\index{Git!fork}
If your own fork carries the \texttt{v0.6} tag, swap \texttt{origin} to
your fork URL on \texttt{dev} with
\texttt{git remote set-url origin git@github.com:<you>/pixelwise.git}
so push and pull continue to target your account.
The HTTPS clone on \texttt{prod} stays read-only.}

\newthought{Verify the public surface.}
From \texttt{dev} or your host browser, hit the HTTPS endpoint and
confirm the frontend renders and the API responds:

\begin{verbatim}
curl -kI https://192.168.56.11
curl -k https://192.168.56.11/api/health
\end{verbatim}

The first call should return \texttt{200 OK} with the security headers
visible; the second should return the health JSON.
You are now at the state where Block~7 ended, with both machines
aligned, ready to start Block~8.

\newthought{Without a snapshot,} replay Block~7 before continuing.
The file state alone is not enough when the next exercises expect Nginx
running with a TLS certificate and the frontend already deployed.


% ==============================================================================================
\newpage
\section{Exercise: CI Workflow}\index{exercises}

\marginnote{\newthought{Exercises} are for practice and reinforcing concepts.
If your VM is broken, restore the \texttt{B7-complete} snapshot.
At the end of this session, take a snapshot named \texttt{B8-complete} as a safety net.}

\newthought{Write the CI workflow.} Create \texttt{.github/workflows/ci.yml} that on every push
lints with \texttt{ruff}\index{ruff} and runs \texttt{pytest} against the boilerplate tests in \texttt{tests/}:

\begin{verbatim}
mkdir -p .github/workflows
nano .github/workflows/ci.yml
\end{verbatim}

Use the YAML from the theory above. Push and watch the workflow run on GitHub.


% ==============================================================================================
\newpage
\section{The Deploy Script}\index{deploy.sh}

\newthought{A standalone deploy script} is reusable from both Actions and manual deploys:

\begin{verbatim}
#!/bin/bash
# scripts/deploy.sh
set -euo pipefail

cd /opt/pixelwise
git pull origin main
source .venv/bin/activate
pip install -r requirements.txt
alembic upgrade head
sudo systemctl restart pixelwise
echo "Deploy complete: $(date)"
\end{verbatim}


% ==============================================================================================
\newpage
\section{Exercise: Deploy Workflow}\index{exercises}

\newthought{Write the deploy workflow.} Create \texttt{.github/workflows/deploy.yml} that on merge
to \texttt{main}, SSHs to the server and runs the deploy script. Add your SSH key and server IP
as GitHub Actions secrets.


% ==============================================================================================
\newpage
\section{Exercise: Deploy Script}\index{exercises}

\newthought{Write the deploy script.} Create \texttt{scripts/deploy.sh} with the steps from the
theory above: pull, install deps, run migrations, restart service.

\begin{verbatim}
mkdir -p scripts
nano scripts/deploy.sh
chmod +x scripts/deploy.sh
\end{verbatim}


% ==============================================================================================
\newpage
\section{Cron: Scheduled Automation}\index{cron}

\newthought{Cron runs commands on a schedule,} no human required. The syntax has five fields:

\begin{verbatim}
# +-------- minute (0-59)
# | +------ hour (0-23)
# | | +---- day of month (1-31)
# | | | +-- month (1-12)
# | | | | + day of week (0-7, Sun=0 or 7)
# | | | | |
  * * * * *  command
\end{verbatim}

\newthought{Common patterns:}
\begin{verbatim}
0 2 * * *    # daily at 2:00 AM
*/5 * * * *  # every 5 minutes
0 0 * * 0    # weekly, Sunday midnight
\end{verbatim}

Edit with \texttt{crontab -e}. List with \texttt{crontab -l}.

\marginnote{Real-world cron uses: database backups, clearing temporary files, health checks that ping your
service and alert on failure, certificate renewal, triggering retraining pipelines, sending summary reports.
The common thread: anything that needs to happen reliably on a schedule without a human remembering.}

\newthought{Students have already seen automation:} \texttt{systemd} restarts their crashed service
automatically (Block~5). \texttt{journalctl} rotates its own logs. Package managers run unattended
security updates. Cron jobs are the explicit, user-controlled version of automation that is already
everywhere under the hood.


% ==============================================================================================
\newpage
\section{Exercise: Nightly Backup Cron}\index{exercises}

\newthought{Set up the nightly backup cron.} On the server, create a cron job that runs nightly
at 2:00~AM. \texttt{pg\_dump} the database, then \texttt{rsync} the dump to the dev VM:

\begin{verbatim}
crontab -e
\end{verbatim}

Add:
\begin{verbatim}
0 2 * * * pg_dump -U pixelwise pixelwise \
    > /tmp/pixelwise_$(date +\%Y\%m\%d).sql \
    && rsync -avz \
    /tmp/pixelwise_$(date +\%Y\%m\%d).sql \
    user@192.168.56.10:/backups/pixelwise/
\end{verbatim}

On the dev VM, create the backup directory:
\begin{verbatim}
mkdir -p /backups/pixelwise
\end{verbatim}


% ==============================================================================================
\newpage
\section{Backup Strategy}\index{backup}\index{pg\_dump}\index{rsync}

\newthought{What to back up:} the database. Code is on GitHub; the model is reproducible from \texttt{train.py}.
The database contains data that cannot be recreated. Every prediction, every timestamp, every confidence score.

\newthought{How:} \texttt{pg\_dump}\index{pg\_dump} exports the database to a SQL file.

\newthought{Where:} off-site. A backup on the same server dies with the server. The \texttt{dev} VM
is the natural off-site target: \texttt{pg\_dump} creates the dump, \texttt{rsync}\index{rsync} copies it
to the dev VM over SSH. The two-machine setup from Block~1 pays off again.

\begin{verbatim}
# Backup script
pg_dump -U pixelwise pixelwise \
    > /tmp/pixelwise_$(date +%Y%m%d).sql
rsync -avz /tmp/pixelwise_$(date +%Y%m%d).sql \
    user@192.168.56.10:/backups/pixelwise/
\end{verbatim}

\marginnote{\texttt{rsync}\index{rsync} (\url{https://rsync.samba.org/}) is efficient file sync over SSH, incremental
(only transfers changed bytes), fast, scriptable. The right tool for moving backup files between machines.}

\newthought{How to verify:} restore from a backup at least once. An untested backup is not a backup:

\begin{verbatim}
# On the dev VM: restore into a test database
psql -U pixelwise -d pixelwise_test \
    < /backups/pixelwise/pixelwise_20260405.sql
\end{verbatim}


% ==============================================================================================
\newpage
\section{Exercise: Verify the Backup}\index{exercises}

\newthought{Verify the backup.} Run the backup manually once. Then restore it into a test database
on the dev VM and verify the data is intact.


% ==============================================================================================
\newpage
\section{Feature Flags}\index{feature flag}

\newthought{Deploying new behaviour without deploying new code.} The simplest form: an env var that
controls which code path runs. Change the flag, restart the service, and the new behaviour is live.
No code change, no PR, no CI pipeline.

\begin{verbatim}
# .env
MODEL_VERSION=v1

# In classifier.py
model_path = f"models/digit_classifier_{MODEL_VERSION}.pkl"
\end{verbatim}

\marginnote{This is how teams roll out features gradually or switch between model versions without risk.
Both \texttt{v1} and \texttt{v2} model files are available. Switching the model is a config change plus
service restart, a hook for comparing model performance later in the course.}


% ==============================================================================================
\newpage
\section{Exercise: Feature Flag for Model Version}\index{exercises}

\newthought{Add a feature flag for model version.} Add \texttt{MODEL\_VERSION=v1} to \texttt{.env}
and \texttt{.env.example}. Update \texttt{classifier.py} to load the model path dynamically.
Place both \texttt{v1.pkl} and \texttt{v2.pkl} in \texttt{models/}. Switch the flag to \texttt{v2},
restart the service, and verify the model changed:

\begin{verbatim}
curl https://192.168.56.11/api/health -k
# Should show "model_version": "v2"
\end{verbatim}


% ==============================================================================================
\newpage
\section{Optional: Webhooks}\index{webhooks}

\newthought{Webhooks are HTTP endpoints triggered by external events.} The internet is event-driven.
GitHub can call your server when code is pushed; Discord can receive a message when your deploy succeeds.
The pattern is always the same: an event happens $\to$ an HTTP POST is sent to a URL you control.


% ==============================================================================================
\newpage
\section{Exercise: Trigger the Pipeline}\index{exercises}

\newthought{Commit and trigger the pipeline.} Push everything you have built in this session and
watch the automation run end to end:

\begin{verbatim}
git add .github/ scripts/ tests/ .env.example
git commit -m "Add CI/CD pipeline, backup cron, feature flag"
git push
\end{verbatim}

Watch the CI workflow on GitHub. If green, merge. The deploy workflow should update the server automatically.

\newthought{Take a snapshot} of both VMs. Name them \texttt{B8-complete}.

\marginnote{\newthought{Where this leaves you.}
The pipeline now runs without you. Every push is linted and tested, every merge deploys, every night
the database is dumped off-site. Automating the system is one half of operating it reliably;
observing it is a thread we will pick up later.}


% ==============================================================================================
\newpage
\section{Security Thread}

\newthought{Store all deployment credentials} (SSH keys, server IP) as GitHub Actions secrets, never in
workflow files. \texttt{rsync} and SSH-based deploys only, never plain protocols. Backup files contain your
entire database. Protect them accordingly with file permissions and consider encryption for off-site storage.


% ==============================================================================================
% BLOCK 3: CONSOLIDATION
% ==============================================================================================

\section{Self-Reflection and Recap}\index{reflection}\index{recap}

\newthought{Reflection questions:}
\begin{itemize}
    \item What is the difference between Continuous Integration and Continuous Deployment?
    \item Why should secrets never appear in workflow YAML files?
    \item Why is a backup on the same server not a real backup?
    \item How do feature flags let you deploy new behaviour without deploying new code?
    \item What happens if the CI pipeline fails? Does the deploy still happen?
\end{itemize}

\newthought{Milestone:} Every push to \texttt{main} automatically lints and tests. Every merge deploys to
the server. Nightly database backup is running and has been verified by restoring once. Feature flag for
model version switching is in place.
